\begin{table}[!htbp]
\centering\small
\caption{Primitive Predictions and Optimized Adjustment Options}
\label{tab:r7_primitive}
\begin{tabular}{rrrrrr}
\toprule
$p$ & $k$ & Lower option & Actual option & Min. shift & Test $d$ \\
\midrule
0.06 & 20 & 0.0261621 & 0.0311074 & 0.0523241 & 0.723838 \\
0.06 & 40 & 0.016882 & 0.0185198 & 0.0337641 & 0.733118 \\
0.06 & 80 & 0.00973227 & 0.0102147 & 0.0194645 & 0.740268 \\
0.08 & 20 & 0.0681238 & 0.0811337 & 0.136248 & 1.43188 \\
0.08 & 40 & 0.0441692 & 0.0488284 & 0.0883384 & 1.45583 \\
0.08 & 80 & 0.0257477 & 0.0271885 & 0.0514953 & 1.47425 \\
0.10 & 20 & 0.12072 & 0.147117 & 0.24144 & 2.12928 \\
0.10 & 40 & 0.0800377 & 0.0900327 & 0.160075 & 2.16996 \\
0.10 & 80 & 0.0475675 & 0.0507768 & 0.095135 & 2.20243 \\
\bottomrule
\end{tabular}
\par\smallskip
\begin{minipage}{0.98\linewidth}\footnotesize
Notes: Two-stage CRRA family: $u_0=2$, $|\theta|\leq0.2$, $C_-=0.5$, $C_+\in\{0.05,0.8\}$, and $A_+-A_-=0.5$. Option denotes $\Omega_+-\Omega_-$. The primitive lower bound determines the test contract, at which every row has $\Delta^0<0<\Delta^{\rm adj}$. The continuous-family proof covers all $p\in[0.06,0.10]$ and $k\in[20,80]$, not just these rows. Values are in fixed cardinal utility and operating-benefit units.
\end{minipage}
\end{table}
